\section{Conclusions}

\begin{frame}{Conclusions}
	Dans ce cours, nous avons fait des rappels de probabilités:
	\begin{itemize}
		\item Probabilités conditionnelles.
		\item Théorème de Bayes.
		\item Facteur de Bayes.
		\item Récursivité du théorème de Bayes.
	\end{itemize}
	\pause
	Nous avons également vu comment l'algorithme Bayésien Naïf utilisait le théorème de Bayes pour résoudre un problème de classification:
	\begin{itemize}
		\item Avec des données suivant une distribution multinomiale (données textuelles): algorithme Bayésien Naïf multinomial.
		\item Avec des données suivant une distribution Gaussienne: algorithme Bayésien Naïf Gaussien.
	\end{itemize}
	\pause
	Nous avons également vu que l'hypothèse de naïveté sur l'indépendance conditionnelle permettait de traiter des problèmes de classification avec des caractéristiques d'entrée suivant des distributions de différentes formes (multinomiale, Gaussienne...) en combinant plusieurs algorithmes Bayésiens Naïfs.
\end{frame}

\begin{frame}{Conclusions}
	Enfin nous avons introduit des techniques rudimentaires de traitements de données textuelles:
	\begin{itemize}
		\item Tokenisation et mots vides.
		\item Lemmatisation et analyse grammaticale.
		\item Vectorisation par comptage et TF-IDF.
	\end{itemize}
	\pause
	Ces traitements ne sont pas exclusifs aux algorithmes Bayésiens Naïfs et pourront être réutilisés pour d'autres algorithmes d'apprentissage automatique.
\end{frame}
